\chapter{Сквозные разборы: как память работает на цельном сценарии}
\label{ch:casestudies}

\section{Зачем нужны сквозные разборы}

Когда архитектура описана по главам, возникает полезная, но опасная иллюзия: будто
субстрат, retrieval, настроение, эпистемика, рантайм и ревизия существуют по
отдельности. В реальном ходе работы памяти они вступают в действие одновременно.
Поэтому здесь полезно пройти несколько сценариев от начала до конца.

\section{Сценарий 1: обещание как факт, обязательство и будущий триггер}

\subsection*{Событие}
В разговоре появляется фраза: «нужно завтра отправить Маше черновик».

\subsection*{Что возникает в памяти}
\begin{enumerate}
\item observation-узел о произнесённом обязательстве;
\item сущность ``Маша'' либо связывается с существующим узлом, либо создаётся;
\item task-like узел ``отправить черновик'';
\item темпоральная привязка ``завтра'';
\item отношение типа \texttt{commitment\_to} и/или \texttt{deadline}.
\end{enumerate}

\subsection*{Почему это не просто эпизод}
Если система запишет только буквальный эпизод разговора, она будет помнить, что
что-то когда-то обсуждалось, но не сможет вовремя действовать. Если же она сразу
запишет только сухую задачу, потеряется provenance и будет трудно понять, откуда
обязательство взялось и можно ли его позже снять.

\begin{principle}[У обязательства минимум два лица]
Для памяти обязательство должно существовать и как событие происхождения, и как
операциональный объект будущего действия.
\end{principle}

\subsection*{Какие слои включаются}
\begin{itemize}
\item write-gate пропускает запись как task-critical;
\item runtime кладёт краткоживущий foreground-след;
\item timescale-слой создаёт будущий временной триггер;
\item retrieval later поднимет задачу не потому, что она семантически похожа на
текущий диалог, а потому что наступает её окно активации.
\end{itemize}

\section{Сценарий 2: ложная гипотеза и каскад ревизии}

\subsection*{Событие}
Рин несколько раз видит человека $Z$ рядом с университетом и рано заключает:
``$Z$ учится здесь''.

\subsection*{Как рождается ошибка}
\begin{enumerate}
\item эпизоды наблюдений проходят write-gate;
\item внутренний walkthrough находит паттерн совместной встречаемости;
\item fusion предлагает гипотезу ``$Z$ связан с университетом'';
\item на слабом материале коммиттер принимает её как \emph{гипотезу с умеренной уверенностью}, а не как факт.
\end{enumerate}

\subsection*{Почему это хороший, а не плохой путь}
Память, которая никогда не рождает таких гипотез, стерильна. Ошибка начинается не в
самом рождении гипотезы, а в преждевременном коллапсе её статуса в ``установленный
факт''.

\subsection*{Контрсвидетельство}
Позже приходит наблюдение: $Z$ не студент, а курьер, часто бывающий в этом районе.

\subsection*{Что должно произойти}
\begin{enumerate}
\item прежняя гипотеза понижается;
\item потомки, опиравшиеся на неё, получают \texttt{need\_to\_revise};
\item provenance сохраняет след того, почему система раньше думала иначе;
\item если таких ошибок много, evidence-aware revision понижает доверие к данному
типу ``соседство $\to$ роль''-выводов.
\end{enumerate}

\section{Сценарий 3: настроение меняет маршрут, но не право на ложь}

\subsection*{Событие}
У Рин понижен фон настроения. Она вспоминает события последних дней и чаще тянется
к неприятным эпизодам.

\subsection*{Нормальная часть этого эффекта}
Mood-congruent recall действительно должен смещать walk. В противном случае система
станет не похожей на живую память, а на бесплотный индекс.

\subsection*{Опасная часть}
Если смещение не ограничено, retrieval начинает собирать контекст почти целиком из
негативно валентных узлов. Тогда:
\begin{enumerate}
\item разговорные ответы становятся систематически мрачнее;
\item метакогниция ошибочно выводит ``у меня всё ухудшается'';
\item возникает риск самосбывающейся гипотезы.
\end{enumerate}

\subsection*{Как работает предохранитель}
\begin{enumerate}
\item настроение смещает walk, но не отменяет epistemic gates;
\item self-facts остаются провизорными, пока не подтверждены независимыми узлами;
\item фоновые detector'ы оценивают автокорреляцию и глубину руминативной петли;
\item при превышении порога включается \textsc{Reflecting}.
\end{enumerate}

\section{Сценарий 4: точное удаление и судьба производных следов}

\subsection*{Событие}
Поступает команда забыть конкретный чувствительный фрагмент недавнего разговора.

\subsection*{Чего недостаточно}
Недостаточно удалить сырой текст или зашифрованный raw-log. Если содержание уже
успело породить узлы, summary, embeddings, pending-candidates и рабочие ассоциации,
память останется загрязнённой производными следами.

\subsection*{Что должна сделать система}
\begin{enumerate}
\item удалить или крипто-шреднуть исходный первичный носитель;
\item найти все производные артефакты по provenance и event lineage;
\item пересчитать affected context indices;
\item проверить, не остались ли осиротевшие потомки без объяснимой опоры.
\end{enumerate}

\begin{pitfall}[Удаление без ревизии осиротевших потомков]
Самая тихая ошибка — не то, что удалённый факт остался лежать в одном месте, а то,
что его производные потомки продолжают жить как будто ничего не произошло. Тогда
формально система ``забыла'', а содержательно — нет.
\end{pitfall}

\section{Сценарий 5: от вопроса пользователя к ответу с provenance}

\subsection*{Вопрос}
``Почему система решила, что этому человеку важна музыка?''

\subsection*{Нормальный плохой ответ}
Система выдаёт гладкую словесную рационализацию без доступа к реальному пути
возникновения убеждения.

\subsection*{Нормальный хороший ответ}
Система может сказать примерно следующее по структуре:
\begin{enumerate}
\item есть два наблюдения;
\item есть один эпизодический паттерн;
\item есть одна осторожная гипотеза;
\item уверенность умеренная;
\item вот что могло бы её поднять или опровергнуть.
\end{enumerate}

\subsection*{Почему это важно}
Такой ответ делает память не только продуктивной, но и проверяемой. Читатель видит
не просто результат, а локальное доказательство, и тем самым память начинает быть
ближе к научной записи, чем к черному ящику.

\section{Вывод}

Сквозные сценарии показывают общий закон курса: память хороша не тогда, когда у
каждого её модуля красивое локальное описание, а тогда, когда между записью,
обходом, ревизией, действием и объяснением нет немых провалов.
